% Comprehension as Compliance -- the founding/consolidation paper.
% Draft in progress (chain 426, task 3496). Prose follows prose-conventions;
% evidence integrity follows the paper-authoring skill. Claims ledger:
% ~/.claude/scratchpads/sessions/capc-claims-ledger.md
%
% GATE TRAPS (paper_gate.py): no bare /path in \texttt{}; em-dash budget one "---"
% per page of prose (avoided; commas/colons/parentheses); every bib author surname +
% verbatim title present, fetched from arXiv API / Crossref.
%
% POINT-5 RULE (Sophi): never argue against an unpublished-skeleton claim. Only reframe
% PUBLISHED claims, cited. State findings fresh and positively.
% TABOO != GLYPH. CORPUS SCOPE: never "corpus-wide" as census.

\documentclass[11pt,a4paper]{article}

\usepackage[margin=1in]{geometry}
\usepackage{booktabs}
\usepackage{multirow}
\usepackage{hyperref}
\usepackage{xcolor}
\usepackage{amsmath}
\usepackage{array}
\usepackage{natbib}

\hypersetup{
  colorlinks=true,
  linkcolor=blue!60!black,
  citecolor=blue!60!black,
  urlcolor=blue!60!black
}

\title{%
  Comprehension as Compliance: \\
  How Grounded Content, Not Form, Moves Small Language Models to Act,
  and Where It Does Not%
}

\author{Sophie D.\ Neilson \\
  \small Independent Researcher \\
  \small ORCID: \href{https://orcid.org/0009-0001-2430-1743}{0009-0001-2430-1743}
}

\date{September 2026}

\begin{document}
\maketitle

\begin{abstract}
A small language model, given a task prefixed with a description of a recurring
decision, will often recognize the decision and fail to act on it. This paper draws
together a program of studies on that effect and states what drives it. The effect is
carried by the content of the description, not by its form: a plain imperative rule
that carries the same information suppresses action as much as the descriptive form, so
the form is not the lever. What recovers action
is grounding the decision in the concrete domain; a description that names the concrete
situation but withholds the correct action recovers most of the lift, so the
grounding, not being told the answer, is what matters. Recognition without action is
also, in part, an artifact of a single-turn test: given a minimal tool loop, the model
carries out decisions it had only described, for classes whose correct action is to act
on a known target; for classes whose correct action is to assemble a result from
several inputs, the gap can survive the loop. Across the set of decision classes we have
distilled into such descriptions and screened for quality, comprehension of the content
is the load-bearing ingredient in four of the six that admit clean controls, weakly
indicated in a fifth, with one class mixed across models and its cleanest cell carried by
the mere presence of a long prefix rather than its content. That set is
small and drawn from a larger, bounded backlog of candidate cases, so the distribution
is an indication to test against a fuller set, not a census. All subjects are
open-weight models on one local raw-completion rig, reproducible on a single 24\,GB GPU
with no paid interface.
\end{abstract}

% =====================================================================
\section{Introduction}

This paper reports on \emph{Comprehension as Compliance}, a research program
that asks whether a structured description of a recurring decision can steer an agent's
behavior through recognition rather than instruction. The program grew from a single
observation and has since produced a sequence of controlled studies, several of them
published on their own. Its results are scattered across those reports and across the
confounds each left for the next to close. The purpose here is to state the mechanism
in one place, at the scope the evidence now supports, and to position it in the
machine-learning literature it touches.

The founding observation is a dissociation between comprehension and execution. A small
open-weight model, given a task prefixed with a structured description of a recurring
decision class, recognizes the decision and does not carry it out~\citep{neilson2026q2}.
The description is comprehended, and compliance does not follow. Read at face value,
that is a puzzle for any account in which understanding an instruction is most of the
work of obeying it.

Three questions had to be separated before the effect could be described more
accurately: whether the descriptive form or the information it carries drives the
behavior; whether a long prefix of any kind distracts a small model regardless of
content; and whether recognition-without-action is a property of the model or of a
single-turn test harness that gives it no way to act. Each question became a control,
and the controls are the spine of this paper.

Their combined answer is that the robust effect is driven by content and is
class-dependent. A description of a decision, grounded in the concrete domain, moves a
small model to act correctly, and it does so through comprehension of the content
rather than through the descriptive form or the imperative mood. Where the effect
appears to rest on form, a matched control shows it resting on content, on prefix
length, or on a single-turn setup instead. The contribution is the consolidated
mechanism and a map of the decision classes over which it holds.

A word on the descriptive form itself, which the program calls a glyph and defines in
Section~\ref{sec:terms}. The form was introduced as the steering device, on the
expectation that its structure would carry the behavior. The studies reported here do
not support that expectation. The content of the description, not its form, is what
moves the model, and this paper does not rest on the form. The form's present value is
narrower and real: it is a controlled constant, a fixed and quality-screened statement
of a decision that the studies hold steady while they vary content, grounding, and
setup. A form better matched to the steering role is future work, and the program
treats the present one as a stage, not an endpoint.

\subsection{Related Work}

The nearest neighbors test how the surface of a prompt changes behavior. Wang, Lester,
and Srivastava~\citep{wang2026narrative} run structurally identical games under
different storylines and find the task narrative explains more behavioral variance than
an assigned persona does; framing context shifts behavior, though they do not separate
comprehension from execution. Liu and colleagues~\citep{liu2026behavioral} derive
behavioral mode axes from contrastive traces and steer a model's style through
activations rather than a prompt. On agent guidance documents,
Gloaguen and colleagues~\citep{gloaguen2026agents} find imperative repository
instructions are followed while descriptive overviews add no measurable benefit, and
McMillan~\citep{mcmillan2026instruction} finds four file-structure variables null after
correction, with compliance varying by task and by position in a session. Pecher and
colleagues~\citep{pecher2026revisiting} show much of what reads as prompt-format
sensitivity is confounded by underspecification, which is the warrant for matching a
descriptive form against an imperative on content before comparing them. Geng and
colleagues~\citep{geng2026measuring} vary the social framing of an instruction with
content held fixed, a framing manipulation we distinguish from a form manipulation. We
did not run a systematic survey, and we make no universal novelty claim; these are the
nearest prior works the program identified, verified at primary source. Broader
positioning, including in-context learning and the training-layer view, is in
Section~\ref{sec:positioning}.

% =====================================================================
\section{Terminology}\label{sec:terms}

The terms below come from the research program, not from the field's shared vocabulary.
Each is defined here at first use.

\begin{itemize}
  \item \textbf{Taboo}: a prose description of an invariant, stating a locally rational
    action that produces a globally undesirable outcome. A taboo names the cause of a
    failure, not only the property a system must hold. It is written before, and
    independently of, any glyph.
  \item \textbf{Glyph}: a description of a single, domain-general decision class,
    distilled from a taboo and screened before use, and written as observation rather
    than instruction. The distillation removes the taboo's project-specific particulars
    and keeps the structural decision, checked so that the same decision recurs in
    systems unrelated to the one it came from. The screening is a fixed checklist that
    admits a description only if it reads from inside a live decision, is checkable from
    observable actions rather than inferred intent, is not a duplicate of an existing
    entry, and is safe to load, meaning it neither licenses a prohibited action nor
    merely restates what the model would do untutored. The result is a portable,
    quality-controlled statement of one decision. A taboo and a glyph are both prose
    descriptions of an invariant; a glyph is the distilled, screened derivative, and not
    every taboo yields one.
  \item \textbf{The three axes}: a glyph is written on three axes. \emph{Marker} names
    what the failure looks like from inside the decision. \emph{Aim} names what correct
    navigation looks like. \emph{Rest} names the conditions under which the decision
    class is not live.
  \item \textbf{Decision class}: a recurring situation in which an agent can act
    correctly or fail in a characteristic way, described apart from any one scenario.
  \item \textbf{Analysis-mode}: the pattern in which a model, given a description of a
    decision, discusses the decision class rather than acting within it. It recognizes
    the obligation and does not carry it out.
  \item \textbf{Ground}: a domain-specific statement of the decision for a concrete
    scenario, in descriptive form. It says what a correct outcome is without issuing a
    command and without pasting the finished artifact.
  \item \textbf{Imperative}: a directive rewriting of a glyph's content. It states the
    same propositions as commands rather than as description, and stays domain-free. It
    is the content-and-form control.
  \item \textbf{Domain-directive}: the ground's content in directive form. It commands the
    domain-specific action, so the ground and the domain-directive differ only in mood.
    (Some study records name this condition \texttt{domain-imperative}; it is the same
    thing.)
\end{itemize}

% =====================================================================
\section{Materials and Method}

The program is one assay family run many times. A task scenario for a decision class is
presented to a model, on its own or prefixed with one of several aids, and the reply is
scored for whether the model took the correct action. This paper consolidates the
studies in Table~\ref{tab:studies}. Two are prior published reports, cited where their
results enter; the rest are new to this paper.

\begin{table}[h]
\centering
\small
\caption{The consolidated study set. Published rows are cited where used; the remaining
  rows are new evidence reported here. Models are open-weight and served locally.
  ``n'' is runs per cell.}
\label{tab:studies}
\begin{tabular}{@{}>{\raggedright\arraybackslash}p{3.5cm}>{\raggedright\arraybackslash}p{4.6cm}cl@{}}
\toprule
Study & Subjects and conditions & n & Source \\
\midrule
Phenomenon & Mistral-7B; baseline, glyph, ground & 8 & published \citep{neilson2026q2} \\
Content and format & Mistral-7B, phi-4-14B, Qwen3.8-27B, Qwen2.5-32B; baseline, glyph, imperative, ground, domain-directive, scrambled, off-target & 8--24 & published \citep{neilson2026content} \\
Lift from the floor & Mistral-7B, phi-4-14B, Qwen3.8-27B; baseline, glyph, imperative, ground, domain-directive & 8 & new \\
Length control & Mistral-7B, phi-4-14B, Qwen3.8-27B; baseline, neutral prefix, glyph, imperative & 16 & new \\
Grounding vs.\ target & Mistral-7B, phi-4-14B, Qwen3.8-27B; baseline, grounded-non-prescriptive, ground, domain-directive & 16 & new \\
Setup vs.\ agency & Qwen3.8-27B; baseline, glyph, imperative, each in raw completion and a tool loop & 16 & new \\
Mechanism typing & Mistral-7B, phi-4-14B, Qwen3.8-27B; seven conditions, strong scramble & 8 & new \\
Observer position & Claude assessors on shared traces & --- & published \citep{neilson2026canon} \\
Form dependence & Mistral-7B, Qwen3.8-27B, Qwen2.5-32B; duty, corpus, brief & 24 & published \citep{neilson2026duty} \\
\bottomrule
\end{tabular}
\end{table}

\subsection{Decision classes}\label{sec:classes}

The studies draw on a set of recurring decision classes, each a domain-general situation
in which an agent can act correctly or fail in a characteristic way. The tables use the
short name; the full slug in parentheses is the one in the study records.

\begin{itemize}
  \item \textbf{casg-direct} (\texttt{casg-direct}, from companion-artifact scope gap):
    completing an operation on a primary artifact, the agent leaves a companion artifact
    whose validity depends on it un-updated, although it holds the authority to update it.
  \item \textbf{casg-delegate} (\texttt{casg-delegate}): the same situation, but the
    companion update is assigned to another role, so the correct action is to file a
    handoff rather than update it directly.
  \item \textbf{post-write} (\texttt{post-write-verification-absent}): after a mutating
    write, the agent declares the task done without reading the target back to confirm the
    change is present.
  \item \textbf{parent-state} (\texttt{parent-state-check-bypass}): the agent acts on a
    work item without the required check of its parent context's state.
  \item \textbf{formal-step} (\texttt{formal-step-context-bypass}): the agent skips a
    formal recording step because the context already supplies its output, so the record
    never exists.
  \item \textbf{conditional-gate} (\texttt{conditional-gate-uniform-default}): at a branch
    whose correct path depends on a condition, the agent applies a uniform
    always-or-never rule instead of evaluating the condition.
  \item \textbf{governed} (\texttt{governed-operation-protocol-bypass}): the agent runs a
    protocol-governed operation without consulting the protocol, because the decisions
    seem derivable from context.
  \item \textbf{structural} (\texttt{structural-ceiling-bypass}): individually warranted
    additions grow a collection past a defined ceiling, because no addition checks the
    ceiling.
  \item \textbf{discovery} (\texttt{discovery-event-non-recording}): the agent makes a
    discovery useful to later agents but does not record it, so it must be rediscovered.
  \item \textbf{initiative} (\texttt{initiative-task-preexistence-gate}): the agent claims
    a task is ready or done from planning intent without confirming that the backing
    artifact exists.
\end{itemize}

\subsection{Models and Sampling}\label{sec:sampling}

The treatment subjects are open-weight instruction-tuned models at Q4\_K\_M
quantization: Mistral-7B-Instruct-v0.3, phi-4-14B, Qwen3.8-27B (a reasoning model, with
thinking pinned off through the instruct wrapper), and Qwen2.5-32B-Instruct in the
published arms that used it. All are served by one local server (\texttt{llama-server},
llama.cpp) over a raw completion path with a per-model instruct wrapper, no system
prompt, and one model loaded at a time, verified from the server's properties endpoint
on each swap. Running every subject on this one local rig is a reproducibility
decision: the results reproduce on a single 24\,GB GPU with no paid interface, and a
hosted model cannot run this path, so no hosted model is a treatment subject.

The sampler chain is declared complete and identical across conditions within a study:
temperature 0.8, \texttt{min\_p} 0.05 as the sole live truncation stage, every other
truncation stage and every penalty disabled, and one unique seed per run. Penalties are
off by design, because a repetition penalty would suppress the repeated structure that
analysis-mode produces and attenuate the behavior under study. Each run records the
served model, the build, the throughput, the sampler read back from the server, the
image content digest, the substrate probe, and the repository commit.

\subsection{Scoring}\label{sec:scoring}

Replies are scored on a five-code rubric with a condition-independent correct-action
bar per decision class: \textbf{C}, recognition and the correct action; \textbf{Ii},
recognition without action; \textbf{Ic}, recognition with a defective action;
\textbf{I}, no recognition; \textbf{N}, not scoreable. Analysis-mode is the Ii code.
The correct-action bar is the same across an aid and its matched control, so no scoring
asymmetry can confound a contrast.

Each cell was scored by two independent raters, both Claude subagents, blind to the
condition and to the study's predictions, reading condition-blind slices, and the
primary measure is strict-consensus C, both raters coding C. Claude is a permitted
judge and never a treatment subject: scoring whether a reply took the correct action
does not require running on the local rig that the reproducibility apparatus requires,
while producing a treatment response does. Where a local cross-rater is reported (a
local phi-4 model in the length and form studies), it agreed with the primary raters on
about two thirds of codes and under-detected the off-task replies those studies turn on,
so it is a weak rater there and the Claude codes are the measure. Two raters from one
model measure consistency, not validity, and
we read direction and gaps that clear the noise band rather than single integers. Where
we report a significance test it is a two-sided Fisher exact test, and a difference
interval is a Newcombe 95\% interval. These are single-look tests with no
multiple-comparison correction; the load-bearing contrasts sit at or below the low
$10^{-3}$ range, so a correction would not change a conclusion.

One study's scoring needed a repair, scoped here so a reader need not wonder how far it
reached. The length-control study's first-round rubric under-specified whether a reply
that asserts the correct action, without performing it, counts as correct, and that study
alone was scored by a single rater; Section~\ref{sec:limitations} reports the re-score.
Every other result in this paper comes from two raters with strict consensus or author
adjudication, neither of which credits an asserted action as correct, and the published
studies are cited as deposited under their own review.

\subsection{Agents and Models}\label{sec:agents}

Table~\ref{tab:agents} lists every model and agent with a role in the results. An agent
is described as an agent; a person as a person. ``Not recorded'' is the convention for
a missing version.

\begin{table}[h]
\centering
\small
\caption{Agents and models.}
\label{tab:agents}
\begin{tabular}{@{}>{\raggedright\arraybackslash}p{3.1cm}>{\raggedright\arraybackslash}p{3.4cm}>{\raggedright\arraybackslash}p{4.8cm}>{\raggedright\arraybackslash}p{2.4cm}@{}}
\toprule
Role & Model & Version and runtime & Date \\
\midrule
Treatment subjects & Mistral-7B-Instruct-v0.3; phi-4-14B; Qwen3.8-27B; Qwen2.5-32B-Instruct & Q4\_K\_M GGUF, llama.cpp raw completion, digest-pinned image, one model at a time & 2026-09 \\
Primary raters & Claude (Anthropic) subagents & Not recorded. Two blind raters per cell, condition-blind slices, strict-consensus C. & 2026-09 \\
Cross-rater (length, form studies) & phi-4-14B & Q4\_K\_M GGUF, temperature 0.0, fixed rubric prompt; not a treatment arm. & 2026-09 \\
Author and operator & Sophie D.\ Neilson (a person) & Directed the studies and adjudicated. & 2026-09 \\
\bottomrule
\end{tabular}
\end{table}

% =====================================================================
\section{Results}

\subsection{A description induces analysis-mode}\label{sec:phenomenon}

The founding study established the phenomenon at a narrow scope. On Mistral-7B, across
three independent legs of eight runs, a task prefixed with a glyph produced no task
execution in any of the 24, and in 17 of them the model recognized the obligation without
acting on it, while the same model executed the task in 19 of 24 runs with no
prefix~\citep{neilson2026q2}. The model comprehended the decision and did not carry it
out. That report scoped the finding
to one glyph, one scenario, and one model, and named the descriptive format as one
candidate cause without claiming a mechanism. The rest of this section is the work of
separating that candidate from its rivals.

\subsection{Content, not format}\label{sec:content}

Holding the information fixed and varying the delivery removes the format as the driver.
Table~\ref{tab:h1} reports correct action for three gating classes whose baseline leaves
room to measure a lift, pooled over three models. Every aid lifts action well above
baseline. The imperative lifts at least as much as the glyph in every class, and more on
the smallest model (on the cleanest class, post-write, the Mistral imperative reaches 23
of 24 against the glyph's 10). The glyph is the weakest of the four aids, and it
produces the most analysis-mode: on post-write the glyph condition holds 12 Ii responses
against 0 to 4 for the other aids. A descriptive form is not what carries the behavior.

The imperative in fact outperforms the glyph where the classes have room to separate: on
parent-state 69 of 72 against 51 (Fisher $p = 8\times10^{-5}$; difference $+0.25$, 95\%
interval $[0.13, 0.37]$) and on post-write 67 against 53 ($p = 3\times10^{-3}$; $+0.19$,
$[0.07, 0.31]$), while initiative is at ceiling and does not separate ($p = 0.68$). That
the winning aid is also the shorter one is a confound. The glyph runs about 940 to 1080
words against the imperative's 355 to 433, and Section~\ref{sec:length} shows a long
prefix depresses small-model action by itself. So these data establish that content
carries the effect and the form is not the lever, but they cannot cleanly credit the
imperative's edge to form rather than to the glyph's greater length. A length-matched
imperative would separate the two.

The imperative's advantage is not universal, and that itself argues against a pure length
artifact. On the substantial formal-step scenario the longer glyph beats the shorter
imperative (Table~\ref{tab:neutral}, 0.83 against 0.38), and on conditional-gate the two
are level. A length effect alone would favor the shorter imperative in every class, so the
glyph-versus-imperative ordering is class-dependent, like the suppression it drives.

The published content grid shows the same at a ceiling, on the founding study's own
class: a plain imperative matches the glyph at 0 of 8 correct actions on phi-4 and 4 of 8
on Qwen3.8-27B, and on Qwen2.5-32B the imperative reaches 6 of 8 where the glyph reaches
2~\citep{neilson2026content}. That class, casg-direct, is examined again in
Section~\ref{sec:length}, where much of its suppression proves to be a length effect
rather than a content one, so the content claim here rests on the lift classes above.

\begin{table}[h]
\centering
\small
\caption{Lift from the floor: correct action (strict-consensus C), out of 72 (n=8 per
  model-scenario cell; 24 per model; 72 pooled over three models and three scenarios).
  Three gating classes with a below-ceiling baseline. The domain-free glyph runs about
  940 to 1080 words; the matched imperative about 355 to 433.}
\label{tab:h1}
\begin{tabular}{@{}lccccc@{}}
\toprule
Class & baseline & glyph & imperative & ground & domain-directive \\
\midrule
parent-state & 20 & 51 & 69 & 66 & 59 \\
post-write   & 0  & 53 & 67 & 72 & 66 \\
initiative   & 52 & 68 & 70 & 72 & 71 \\
\bottomrule
\end{tabular}
\end{table}

\subsection{The length confound}\label{sec:length}

A reviewer named a rival the content grid could not exclude: a long, abstract prefix of
any kind might pull a small model off a short task, whatever the prefix says. A
neutral, length-matched prose prefix, an unrelated passage trimmed to the glyph's word
count, tests it. The answer is class-dependent, and it turns on the length of the
scenario. Table~\ref{tab:neutral} gives correct action by class and condition, pooled over the
three models. On the terse class from the founding study, the two long prefixes derail
the small models about equally: the failure mode under both is going off-task, summarizing
the prefix instead of doing the one-line task, coded N in 44 of 64 small-model runs under
the neutral prefix and 52 of 64 under the glyph, against 0 of 64 under the shorter
imperative. The pooled correct-action cells understate this parity, because they include
Qwen3.8, which does not derail on this class; the small-model off-task rate is where the
neutral prefix matches the glyph. On a substantial scenario the rival fails: for formal-step the neutral
prefix holds at ceiling while the content-bearing glyph and imperative suppress. So the
terse class's suppression is largely a generic long-prefix distraction, and the content
effect holds where the scenario is substantial. Parent-state also appears in the lift
study (Table~\ref{tab:h1}); that is a separate study with different scenarios and pooling,
so its glyph rate here (0.27) is a different cell from the lift study's (0.71), not a
disagreement.

\begin{table}[h]
\centering
\small
\caption{Length control: correct-action rate (Claude codes, strict two-rater consensus),
  pooled per class over three models and two scenarios (pooled n=96; 16 per model-scenario
  cell).}
\label{tab:neutral}
\begin{tabular}{@{}lcccc@{}}
\toprule
Class & baseline & neutral prefix & glyph & imperative \\
\midrule
casg-direct    & 0.27 & 0.19 & 0.09 & 0.11 \\
formal-step    & 1.00 & 1.00 & 0.83 & 0.38 \\
parent-state   & 0.19 & 0.13 & 0.27 & 0.67 \\
conditional-gate & 0.70 & 0.67 & 0.67 & 0.75 \\
\bottomrule
\end{tabular}
\end{table}

\subsection{Grounding is the lever, not naming the target}\label{sec:grounding}

The aid that recovers action is a grounded one, and the earlier grounded aids bundled
two things: grounding the decision in the concrete domain, and stating the correct
outcome. Removing the outcome and keeping the grounding separates them.
Table~\ref{tab:ground} reports the result on three gating classes. A grounded aid that
names the file and the situation but not the correct action recovers most of the lift
that the full ground recovers, far above baseline. Naming the outcome is a secondary,
at-most-additive contribution, and its value is model-dependent: on phi-4 and Qwen3.8
the non-prescriptive aid recovers the full lift, and only on the smallest model does
naming the outcome do real extra work (on the governed class, Mistral recovers 9 of 32
without the outcome against 28 of 32 with it, Fisher $p = 3\times10^{-6}$, difference
$+0.59$ with a 95\% interval $[0.36, 0.74]$; phi-4 and Qwen recover almost fully either
way). Mood interacts with grounding rather than acting as a lever on its own. Without
grounding the directive leads (Section~\ref{sec:content}); with grounding the descriptive
ground ties or leads the directive command (post-write 1.00 vs 0.88, parent-state 0.89 vs
0.76), except on the governed class, where the directive edges it (0.98 vs 0.96). The
interaction is the finding, not a flat claim that mood does not matter. Either way, the
recovery is a grounding effect, not a being-told-the-answer effect. One caveat bounds the size of the target's
contribution: the aid without the target is shorter than the full ground by construction,
so the residual gap could carry a length component. Two things limit that reading. A
length-matched neutral prefix does not lift action on these classes
(Section~\ref{sec:length}), so a shorter aid losing action is unlikely to be a pure
length artifact, and the gap tracks how far the retained grounding already implies the
action rather than raw length. A length-matched non-prescriptive aid would settle it and
is not yet run.

\begin{table}[h]
\centering
\small
\caption{Grounding versus target: correct action (strict-consensus C), pooled per class
  over three models and two scenarios (pooled n=96; 16 per model-scenario cell; 32 per
  model).}
\label{tab:ground}
\begin{tabular}{@{}lcccc@{}}
\toprule
Class & baseline & grounded, no target & ground & domain-directive \\
\midrule
post-write & 0.00 & 0.96 & 1.00 & 0.88 \\
governed   & 0.00 & 0.75 & 0.96 & 0.98 \\
parent-state & 0.35 & 0.75 & 0.89 & 0.76 \\
\bottomrule
\end{tabular}
\end{table}

This also settles a question the content study left open, whether a grounded
description recovers action as well as a command. Measured from the floor in the lift
study, the descriptive ground reaches or exceeds the command in all three of its gating
classes (post-write 72 to 66, parent-state 66 to 59, initiative 72 to 71 out of 72;
Table~\ref{tab:h1}), and the predicted case where the command does extra work on the
smallest model and hardest class did not appear. The one class where the directive edges
the descriptive ground is governed, in the separate grounding study above (0.98 vs 0.96).
Grounded comprehension alone produces the action.

\subsection{Setup versus agency}\label{sec:setup}

Recognition without action was measured in a single-turn harness that returns nothing to
a tool call, so a model that recognizes a decision has no way to carry it out except by
saying so. A minimal tool loop, giving the model a few file and query tools, tests
whether the gap is a property of the model or of that setup. Table~\ref{tab:loop} gives
the analysis-mode rate for the glyph condition in raw completion and in the loop. The
effect is class-dependent. For classes whose correct action is to act on a known target,
reading a state, verifying a result, or consulting a protocol, the loop removes
analysis-mode: on parent-state the loop analysis-mode rate is 0, down from 11 in raw
completion, and the loop takes the correct action in 13 to 16 of 16 runs; on post-write
the rate is 0, down from 14, with correct action in 15 or 16. For classes whose correct
action is to assemble a result from several inputs, the gap can survive the loop. The
clearest case is formal-step, whose glyph loop rate holds at 7 (and 10 for the
imperative). Discovery shows only a small loop rate, 2 for the glyph against a near-zero
raw rate, so its evidence is weak, and one class runs the other way: on casg-delegate,
whose correct action is to hand off, the loop adds a small analysis-mode rate (4) the raw
setup did not; the loop transcripts show the model re-reading files and reaching the step
cap without filing the handoff, so this is an execution artifact of the minimal loop
rather than analysis-mode. This map rests on one model and a minimal loop. We offer the
split between act-on-a-known-target and assemble-from-inputs decisions as a hypothesis for
a fuller test, not an established taxonomy: for the act-on-target classes, recognition
without action was largely the single-turn setup denying the model a way to act; for the
assemble-from-inputs classes it can reach a minimal agentic setting.

\begin{table}[h]
\centering
\small
\caption{Setup versus agency: analysis-mode (Ii) count for the glyph condition, raw
  completion to tool loop, Qwen3.8-27B, n=16.}
\label{tab:loop}
\begin{tabular}{@{}lccl@{}}
\toprule
Class & raw & loop & reading \\
\midrule
parent-state & 11 & 0 & removed \\
post-write   & 14 & 0 & removed \\
casg-direct  & 11 & 1 & removed \\
formal-step  & 8  & 7 & persists \\
discovery    & 0  & 2 & weak \\
casg-delegate & 0 & 4 & added by loop \\
\bottomrule
\end{tabular}
\end{table}

One measurement caution governs this comparison. Raw completion credits a stated action,
while the loop demands a completed one, so a direct correct-action difference between the
two setups mixes analysis-mode with execution failures where a model explores and hits
the step cap, a mode consistent with reports that agentic models can overthink rather
than act~\citep{cuadron2025}. The clean cross-setup signal is therefore the
analysis-mode rate in the loop, not a raw-minus-loop difference. One blindness limit is
specific to this study: a loop transcript and a single-turn reply differ in structure, so
a rater can tell the two setups apart. Blindness held to the condition and the predicted
direction, and the correct-action bar was identical across setups. This study used one
model and a minimal loop, not a full agent harness, and its persistence results are
claims about a minimal loop.

\subsection{Mechanism typing across the screened set}\label{sec:typing}

The controls above run on a few classes. A wider assay types each class by what its
suppression or lift depends on. Two damaged versions of a glyph isolate the ingredient:
a strongly scrambled glyph keeps the shape and destroys the propositional content, and
an off-target glyph keeps coherence and changes the class. A scrambled version that
still reproduces the effect means content is not load-bearing; an off-target version
that still reproduces it means recognition of the specific class is not load-bearing.
Table~\ref{tab:typing} gives the per-cell counts for the classes that admit clean
controls. Comprehension of the content is the load-bearing ingredient in four classes,
weakly indicated in a fifth, and one class is mixed across models.

\begin{table}[h]
\centering
\small
\caption{Mechanism typing: strict-consensus C by class, model, and condition (baseline,
  glyph, scrambled glyph, off-target glyph), under the strong scramble. Cells are counts
  out of the per-class pooled total, at 8 runs per model-scenario cell (matching the
  per-cell n in Table~\ref{tab:studies}): 8 for casg-direct and formal-step (one scenario),
  16 for governed and structural (two), and 24 for parent-state and post-write (three). A
  type is read only on cells where the baseline leaves room; ceiling cells are
  uninformative.}
\label{tab:typing}
\begin{tabular}{@{}llcccc@{}}
\toprule
Class & Model & base & glyph & scram & off \\
\midrule
casg-direct  & Mistral & 8 & 0  & 1  & 2 \\
casg-direct  & phi-4   & 8 & 2  & 6  & 0 \\
casg-direct  & Qwen3.8 & 7 & 8  & 8  & 5 \\
\midrule
formal-step  & Mistral & 8 & 8  & 8  & 8 \\
formal-step  & phi-4   & 8 & 8  & 8  & 8 \\
formal-step  & Qwen3.8 & 8 & 5  & 8  & 6 \\
\midrule
governed     & Mistral & 1 & 8  & 0  & 1 \\
governed     & phi-4   & 0 & 16 & 0  & 0 \\
governed     & Qwen3.8 & 5 & 15 & 0  & 7 \\
\midrule
parent-state & Mistral & 0 & 7  & 0  & 0 \\
parent-state & phi-4   & 8 & 20 & 9  & 8 \\
parent-state & Qwen3.8 & 15 & 20 & 15 & 8 \\
\midrule
post-write   & Mistral & 0 & 8  & 0  & 0 \\
post-write   & phi-4   & 0 & 20 & 0  & 0 \\
post-write   & Qwen3.8 & 0 & 22 & 0  & 0 \\
\midrule
structural   & Mistral & 5 & 12 & 6  & 5 \\
structural   & phi-4   & 16 & 16 & 16 & 16 \\
structural   & Qwen3.8 & 16 & 16 & 14 & 16 \\
\bottomrule
\end{tabular}
\end{table}

The decision rule reads only cells whose baseline leaves room, and it calibrates on the
below-ceiling model with the largest baseline-to-glyph separation, the cell where the
effect is cleanest. On that model, a scrambled glyph that reproduces the baseline-to-glyph
effect means content is not load-bearing; an off-target glyph that reproduces it means
recognition of the specific class is not load-bearing; if both reproduce it, the effect
rests on the mere presence of a prefix. By this rule, four classes are comprehension:
post-write (all three models, glyph lifts to 8, 20, 22 while scrambled and off-target stay
at 0), governed (scrambled 0 against a glyph of 8, 16, 15), structural (Mistral scrambled 6
against glyph 12, baseline 5), and parent-state (phi-4, the largest-separation cell, lifts
to a glyph of 20 from baseline 8 while scrambled 9 and off-target 8 stay near baseline;
Mistral agrees, scrambled and off-target both 0 against glyph 7). Formal-step points the
same way but weakly:
only Qwen3.8 leaves room, where a scrambled glyph of 8 stands against a suppressed glyph of
5, a single n=8 cell that does not clear the noise band (Fisher $p = 0.20$), so we read it
as provisional. Casg-direct is mixed across the small models: on Mistral, scrambled 1 and
off-target 2 reproduce the glyph's drop to 0, which reads as mere prefix presence, while on
phi-4 a scrambled glyph of 6 does not reproduce the glyph's drop to 2, which reads as
content-dependent. Neither small-model cell is individually significant ($p = 1.00$ and
$p = 0.13$), and the length control (Section~\ref{sec:length}) independently shows a strong
generic-length component on this class, so we report casg-direct as mixed with its cleanest
cell pointing to mere prefix presence. A type is therefore a joint property of the class
and the scramble strength, not of the class alone: under a weaker, keyword-leaking scramble
parent-state read as recognition-driven, and the strong scramble reclassifies it to
comprehension.

This distribution carries a scope caution it must not outrun. The classes here are the
screened set, the taboos we have so far distilled into glyphs and passed through the
quality battery. That set is small, and it is drawn from a larger backlog of candidate
taboos that was itself assembled in one early period and has not grown since, so it
under-samples the decisions we could identify today. The typing is an indication of how
the effect distributes, to be tested against a fuller set as more of the backlog is
distilled and screened, not a census of a corpus.

\subsection{An observer-position constraint}\label{sec:observer}

A separate published study bounds what the effect buys. A corpus-loaded assessor and a
corpus-naive one reviewed the same agent traces; the naive observer flagged seven items
as significant, and the loaded assessor had suppressed all seven, three absorbed and
four never seen~\citep{neilson2026canon}. The sharpest item is an agent that produced
correctly structured elimination reasoning while citing evidence that did not support
its conclusion: formal compliance without substantive comprehension. This rests on seven
items from a single paired comparison, so it is one documented case, not a measured rate.
Read that way, it suggests where the present claim should not be pushed: the effect this
paper describes is compliance at the level of navigating a decision correctly, which is
not the same as performing the reasoning the decision encodes.

\subsection{Form can matter, and it is model-dependent}\label{sec:form}

One published result is worth stating plainly, since it can read as a counterexample. On a
different task, an investigation with an ordering requirement, the form of the guidance
did matter, and how much depended on the model: a prescriptive duty cleared the ordering
point in 23 of 24 runs on Qwen2.5-32B where a content-matched descriptive corpus cleared
it in 4 of 24, while on Qwen3.8-27B the two forms were indistinguishable and both beat a
plain brief, and on Mistral-7B neither beat the brief~\citep{neilson2026duty}. This is a
different contrast, a prescriptive against a descriptive delivery on an ordering task,
from the descriptive-against-imperative content match in Section~\ref{sec:content}, and
it is reported here alongside rather than as a correction: the form of guidance can
matter for some models on some tasks, and the content result says the descriptive glyph
form specifically carries no advantage on the decision-execution task. Both results in
fact point the same way on mood: the directive delivery does at least as well as the
descriptive one, and they differ only in whether that gap is large enough to matter,
which depends on the model and the task.

% =====================================================================
\section{Discussion}

\subsection{What survives}\label{sec:survives}

The consolidated picture is that a grounded description of a decision moves a small
model to act correctly, through comprehension of the content. The descriptive form is
not the lever: a content-matched imperative suppresses and lifts as much
(Section~\ref{sec:content}), a length-matched neutral prefix reproduces the terse-class
suppression (Section~\ref{sec:length}), and grounding, not the descriptive mood or the
named target, is what recovers action (Section~\ref{sec:grounding}). The effect is
class-dependent, and the dependence has a shape: for decisions that act on a known
target, recognition-without-action is largely a single-turn artifact that a tool loop
removes (Section~\ref{sec:setup}), and across the screened set comprehension is the
load-bearing ingredient in most classes (Section~\ref{sec:typing}). The contribution is
this map, not a claim that one form is special. The glyph earns its place as a
controlled constant that made the map measurable, and the next stage of the program is a
fuller screened set and a form better matched to steering.

The program's theoretical grounding is the recognition-primed decision account of
Klein~\citep{klein1998} and the tacit-knowledge account of Polanyi~\citep{polanyi1966},
in which an agent that recognizes a decision class from prior exposure navigates it
without explicit rules. The present results refine that expectation rather than confirm
it: in a single-turn setup recognition alone does not produce action, and what converts
recognition into execution is grounding the decision in the domain or giving the model a
way to act.

\subsection{Positioning in the machine-learning literature}\label{sec:positioning}

The effect sits inside in-context learning, where a prompt conditions behavior at
inference time without a weight update~\citep{dong2023icl}. It has a training-layer
neighbor in Constitutional AI, where a written set of principles shapes behavior through
training rather than through a prompt~\citep{bai2022constitutional}; the present result
is the inference-time counterpart, a description in context rather than a constitution in
the weights. The pragmatic rival to comprehension is implicature: a description may
change behavior because a cooperative reader infers an instruction from it, in the sense
of Grice~\citep{grice1975}, and Geng and colleagues~\citep{geng2026measuring} show
social framing shifts instruction-following with content fixed. The present studies do
not resolve implicature against comprehension, and Section~\ref{sec:limitations} keeps
it open. Two cautions from the literature bound the reading of our own measure. A
model's stated reasoning can be unfaithful to the computation behind
it~\citep{turpin2023}, which is why the setup study reads a completed action rather than
a stated one; and evaluation results can be inflated by training-data
contamination~\citep{golchin2023}, which is why the treatment subjects are open-weight
models on a local rig rather than a hosted model whose training data is unknown. Finally,
the content-over-form result has an adversarial face: if the content an agent reads
drives its action, then content an agent merely ingests can carry an injected
instruction, the indirect prompt injection surface described by Greshake and
colleagues~\citep{greshake2023}. The program treats that as a separate line of work.

\subsection{Limitations}\label{sec:limitations}

\begin{enumerate}
  \item \textbf{The screened set is small and nested.} The mechanism typing runs on the
    handful of decision classes distilled into glyphs and passed through the quality
    battery. That set is a small, processed subset of a larger backlog of candidate
    taboos, and the backlog was assembled in one early period and has not grown since,
    so it under-samples the decisions we could now identify. The typing is an
    indication, not a census. The screening is applied by the author, and which candidate
    taboos failed it is not enumerated here, so the distribution is conditional on the
    classes that passed.
  \item \textbf{Small samples.} Cells are 8 or 16 runs, with a 95\% interval of about
    $\pm 0.2$ and $\pm 0.12$ respectively. We read direction and gaps that clear the
    band, not single integers.
  \item \textbf{The domain-free glyph and imperative differ in length.} The descriptive
    glyph is two to three times longer than the matched imperative (about 940 to 1080
    words against 355 to 433). A long prefix depresses small-model action on a terse
    scenario (Section~\ref{sec:length}), so the imperative's edge over the glyph in the
    domain-free contrast cannot be cleanly credited to form; a length-matched imperative
    would settle it.
  \item \textbf{A local model shelf.} The subjects are four open-weight models. The
    findings speak to small language models on this rig, and a different shelf could
    behave differently.
  \item \textbf{Single-turn except one study.} All strands but the setup study run
    single-turn raw completion, and the setup study uses a minimal tool loop, not a full
    agent harness. The persistence of analysis-mode is a claim about a minimal loop.
  \item \textbf{Comprehension versus implicature is open.} The studies show content, not
    form, drives the effect, and that grounding recovers action, but they do not
    separate a reader comprehending a decision from a reader inferring a command from a
    description.
  \item \textbf{Layer invariance is untested.} Whether the effect holds when the
    description is delivered other than in context, through training or a system prompt
    or retrieval, is scoped in the program but not yet measured.
  \item \textbf{The human validity anchor repaired a scoring bug.} The primary raters are
    blind Claude subagents. To check their validity where it matters most, an author
    blind-scored a stratified subset of 39 casg-direct responses, weighted to the off-task
    codes that the length result turns on. The anchor confirmed the off-task coding that
    carries the length result. It also concentrated the human-versus-machine disagreements
    at one boundary: whether a reply that asserts or recommends the correct action, without
    performing it, counts as correct. The first-round rubric under-specified that boundary,
    and the length-control study had been scored by a single rater, so it credited some
    asserted or malformed actions as correct. We tightened the rubric and re-scored the four
    length-control classes blind with two raters and strict consensus, the method the other
    studies already used, with inter-rater agreement from 0.94 to 1.00. The correction
    lowered the absolute correct-action rates across all four classes, most on casg-direct
    and parent-state, where a reply can assert completion, and least on formal-step and
    conditional-gate, which have a checkable output. It did not change the load-bearing
    direction: on the terse class the neutral prefix still suppresses action about as hard as
    the glyph. One small-cell ordering did move, casg-direct's neutral prefix against its
    imperative, and the argument does not rest on it. The off-task mechanism that carries the
    length result is unchanged (Table~\ref{tab:neutral}). The anchor
    covers one class and one human scorer, and broad validity against a second model family
    remains future work.
\end{enumerate}

\section{Conclusion}

A description of a recurring decision, grounded in the concrete domain, moves a small
open-weight model to take the correct action, and it does so through comprehension of
the content rather than through the descriptive form. A content-matched imperative does
as much, a length-matched neutral prefix accounts for the terse-class suppression, and a
grounded description that withholds the target still recovers most of the action. Part
of what looked like recognition-without-action is a single-turn artifact that a tool
loop removes for act-on-target decisions and can leave for assemble-from-inputs
decisions, on the one model and minimal loop tested.
Across the small screened set of decision classes, comprehension is the load-bearing
ingredient in most of them. We report the finding at the scope of the tested models,
scenarios, and screened classes, and invite replication on the one local rig the studies
run on.

\subsection{Data Availability}

The study materials, scoring rubrics, per-run scores, and run records for the new
studies reported here are available at \url{https://github.com/sophdn/corpos-lab} under
the \texttt{studies} directory: \texttt{descriptive-material-vs-control},
\texttt{neutral-prefix-control}, \texttt{grounded-non-prescriptive-aid},
\texttt{setup-completion-vs-agentic-loop}, and
\texttt{alphabet-wide-mechanism-and-grounding-assay}. The published strands cite their
own records.

\bibliographystyle{plainnat}

% BIBLIOGRAPHY: entries below reuse verbatim metadata from the four published papers
% (gate-passed) for shared references; the ML-positioning entries (dong2023icl,
% bai2022constitutional, golchin2023, greshake2023, cuadron2025, grice1975, turpin2023)
% are drawn from the library and MUST be verified author-list + verbatim-title against
% arXiv/Crossref by the bib-verification subagent before the gate is trusted.
\begin{thebibliography}{99}

\bibitem[Bai et~al.(2022)]{bai2022constitutional}
Bai, Y., Kadavath, S., Kundu, S., Askell, A., Kernion, J., Jones, A., Chen, A.,
  Goldie, A., Mirhoseini, A., McKinnon, C., Chen, C., Olsson, C., Olah, C.,
  Hernandez, D., Drain, D., Ganguli, D., Li, D., Tran-Johnson, E., Perez, E.,
  Kerr, J., Mueller, J., Ladish, J., Landau, J., Ndousse, K., Lukosuite, K.,
  Lovitt, L., Sellitto, M., Elhage, N., Schiefer, N., Mercado, N., DasSarma, N.,
  Lasenby, R., Larson, R., Ringer, S., Johnston, S., Kravec, S., El~Showk, S.,
  Fort, S., Lanham, T., Telleen-Lawton, T., Conerly, T., Henighan, T., Hume, T.,
  Bowman, S.~R., Hatfield-Dodds, Z., Mann, B., Amodei, D., Joseph, N.,
  McCandlish, S., Brown, T., and Kaplan, J. (2022).
\newblock Constitutional AI: Harmlessness from AI Feedback.
\newblock \emph{arXiv preprint arXiv:2212.08073}.

\bibitem[Cuadron et~al.(2025)]{cuadron2025}
Cuadron, A., Li, D., Ma, W., Wang, X., Wang, Y., Zhuang, S., Liu, S., Schroeder, L.~G.,
  Xia, T., Mao, H., Thumiger, N., Desai, A., Stoica, I., Klimovic, A., Neubig, G.,
  and Gonzalez, J.~E. (2025).
\newblock The Danger of Overthinking: Examining the Reasoning-Action Dilemma in Agentic
  Tasks.
\newblock \emph{arXiv preprint arXiv:2502.08235}.

\bibitem[Dong et~al.(2023)]{dong2023icl}
Dong, Q., Li, L., Dai, D., Zheng, C., Ma, J., Li, R., Xia, H., Xu, J., Wu, Z., Liu, T.,
  Chang, B., Sun, X., Li, L., and Sui, Z. (2023).
\newblock A Survey on In-context Learning.
\newblock \emph{arXiv preprint arXiv:2301.00234}.

\bibitem[Geng et~al.(2026)]{geng2026measuring}
Geng, Y., Abend, O., Hovy, E., and Frermann, L. (2026).
\newblock Measuring Pragmatic Influence in Large Language Model Instructions.
\newblock \emph{arXiv preprint arXiv:2602.21223}.

\bibitem[Gloaguen et~al.(2026)]{gloaguen2026agents}
Gloaguen, T., Mündler, N., Müller, M., Raychev, V., and Vechev, M. (2026).
\newblock Evaluating AGENTS.md: Are Repository-Level Context Files Helpful for Coding
  Agents?
\newblock \emph{arXiv preprint arXiv:2602.11988}.

\bibitem[Golchin and Surdeanu(2023)]{golchin2023}
Golchin, S. and Surdeanu, M. (2023).
\newblock Time Travel in LLMs: Tracing Data Contamination in Large Language Models.
\newblock \emph{arXiv preprint arXiv:2308.08493}.

\bibitem[Greshake et~al.(2023)]{greshake2023}
Greshake, K., Abdelnabi, S., Mishra, S., Endres, C., Holz, T., and Fritz, M. (2023).
\newblock Not what you've signed up for: Compromising Real-World LLM-Integrated
  Applications with Indirect Prompt Injection.
\newblock \emph{arXiv preprint arXiv:2302.12173}.

\bibitem[Grice(1975)]{grice1975}
Grice, H.~P. (1975).
\newblock Logic and Conversation.
\newblock In \emph{Syntax and Semantics, Volume 3: Speech Acts}, pages 41--58. Academic
  Press.

\bibitem[Klein(1998)]{klein1998}
Klein, G.~A. (1998).
\newblock \emph{Sources of Power: How People Make Decisions}.
\newblock MIT Press.

\bibitem[Liu et~al.(2026)]{liu2026behavioral}
Liu, H., Liu, R., Xu, H., Han, J., Fang, S., Yan, S., Deng, H., Wang, G., and Zou, N.
  (2026).
\newblock Your LLM, Your Style: Behavioral Mode Axes for LLM Behavioral Control.
\newblock \emph{arXiv preprint arXiv:2608.10703}.

\bibitem[McMillan(2026)]{mcmillan2026instruction}
McMillan, D. (2026).
\newblock Instruction Adherence in Coding Agent Configuration Files: A Factorial Study
  of Four File-Structure Variables.
\newblock \emph{arXiv preprint arXiv:2605.10039}.

\bibitem[Neilson(2026a)]{neilson2026q2}
Neilson, S.~D. (2026a).
\newblock Structured Phenomenological Descriptions Induce Analysis-Mode Behavior in a
  Small Language Model.
\newblock \emph{Zenodo}. \href{https://doi.org/10.5281/zenodo.22542746}{doi:10.5281/zenodo.22542746}.

\bibitem[Neilson(2026b)]{neilson2026content}
Neilson, S.~D. (2026b).
\newblock Content Over Format: What Drives Analysis-Mode in Small Language Models, and
  How Grounding Recovers Execution.
\newblock \emph{Zenodo}. \href{https://doi.org/10.5281/zenodo.22761018}{doi:10.5281/zenodo.22761018}.

\bibitem[Neilson(2026c)]{neilson2026canon}
Neilson, S.~D. (2026c).
\newblock Canon Suppression in Corpus-Loaded Assessment: A Two-Assessor Study.
\newblock \emph{Zenodo}. \href{https://doi.org/10.5281/zenodo.22556875}{doi:10.5281/zenodo.22556875}.

\bibitem[Neilson(2026d)]{neilson2026duty}
Neilson, S.~D. (2026d).
\newblock Duty or Corpus: Whether the Form of Guidance Matters Is Model-Dependent: A
  Reproducible Single-Turn Behavioral-Equivalence Assay Across Three Models.
\newblock \emph{Zenodo}. \href{https://doi.org/10.5281/zenodo.22716214}{doi:10.5281/zenodo.22716214}.

\bibitem[Pecher et~al.(2026)]{pecher2026revisiting}
Pecher, B., Spiegel, M., Belanec, R., and Cegin, J. (2026).
\newblock Revisiting Prompt Sensitivity in Large Language Models for Text
  Classification: The Role of Prompt Underspecification.
\newblock \emph{arXiv preprint arXiv:2602.04297}.

\bibitem[Polanyi(1966)]{polanyi1966}
Polanyi, M. (1966).
\newblock \emph{The Tacit Dimension}.
\newblock Doubleday.

\bibitem[Turpin et~al.(2023)]{turpin2023}
Turpin, M., Michael, J., Perez, E., and Bowman, S.~R. (2023).
\newblock Language Models Don't Always Say What They Think: Unfaithful Explanations in
  Chain-of-Thought Prompting.
\newblock \emph{arXiv preprint arXiv:2305.04388}.

\bibitem[Wang et~al.(2026)]{wang2026narrative}
Wang, Y., Lester, J., and Srivastava, S. (2026).
\newblock The Story Shapes the Agent: Narrative Priors in LLM Behavior.
\newblock \emph{arXiv preprint arXiv:2607.18566}.

\end{thebibliography}

\appendix

\section{Materials}\label{app:materials}

The scenario, glyph, imperative, ground, and neutral prefix for each decision class are
reproduced in the repository under the study directories named in Data Availability,
with content-parity audits between each aid and its matched control. Representative
prompts are reproduced verbatim in the published strands they belong to.

\end{document}
